\runningthread{信源统计特性}

\begin{document}

\loomcover
  {信源编码的统计基础}
  {通信原理：从概率模型到压缩极限}
  {fanzuituanhuo}
  {\today}

\section*{总览地图}
\warmth{0}\quad\whisper{先抓住一条主线：信源编码先建模不确定性，再讨论在什么失真约束下能压到多低。}

\begin{strand}
本页笔记的 spine 是：\keyword{先建立信源统计模型，再选择失真约束，最后用信息度量给出可达到的压缩极限}。
单符号、序列、连续和有记忆信源给出“要编码的对象”；熵、条件熵、联合熵和互信息给出
“能压缩到哪里”的语言。后面的无失真编码、限失真编码和变换编码都会沿着这条线展开。
\end{strand}

\block{一眼对照}
\noindent\begin{tabularx}{\linewidth}{@{}l l L@{}}
\toprule
{\color{indigo}对象} & {\color{indigo}符号} & \multicolumn{1}{l}{\color{indigo}统计描述}\\
\midrule
离散单符号信源 & \(X,\ \Xcal=\{x_1,x_2,\cdots,x_n\}\) & 用概率质量函数 \(\pmf(x_i)\) 描述每个符号出现的概率。\\
连续单符号信源 & \(X\in(a,b)\) & 用概率密度函数 \(\pmf(x)\) 描述取值附近的概率强度（本章不展开）。\\
离散消息序列信源 & \(X^L\) & 用 \(L\) 维联合分布 \(\pmf(x_1,x_2,\cdots,x_L)\) 描述整段序列。\\
无记忆序列信源 & \(X_1,\cdots,X_L\) 相互独立 & 联合概率可以分解为各符号概率的乘积。\\
信息熵 & \(H(X)\) & 度量一个信源平均不确定性，等于自信息量的数学期望。\\
条件熵与联合熵 & \(H(Y|X),\ H(X,Y)\) & 描述已知一个信源后另一个信源的剩余不确定性，以及两个信源整体的不确定性。\\
互信息 & \(\info(X;Y)\) & 描述观察 \(Y\) 后关于 \(X\) 的不确定性减少量。\\
\bottomrule
\end{tabularx}

\block{后续路线}
\noindent\begin{tabularx}{\linewidth}{@{}l l L@{}}
\toprule
{\color{indigo}编码问题} & {\color{indigo}核心量} & \multicolumn{1}{l}{\color{indigo}本章前置知识的作用}\\
\midrule
无失真离散信源编码 & \(H(X)\) & 熵给出平均码长的理论下界，编码要尽量逼近这个下界。\\
限失真信源编码 & \(R(D)\) & 在允许平均失真 \(D\) 时，用率失真函数描述最小所需码率。\\
连续信源限失真编码 & \(R(D)\), \(d(x,\hat x)\) & 连续取值需要指定失真度量和重构变量 \(\hat X\)。\\
矢量量化 & 码本 \(\{y_i\}\), \(D\) & 对多个样值联合量化，利用样值间相关性，优于标量量化。\\
有记忆信源编码 & 相关性、变换 & 先削弱样本间相关性，再对近似独立的分量编码，例如 transform coding。\\
\bottomrule
\end{tabularx}

\section{通信系统与信源编码}
\warmth{0}\quad\whisper{这部分回答：为什么一上来就要研究信源的统计特性。}

\begin{strand}
通信的起点是\keyword{信源}，它产生待传输的信息。信道只是承载信号的通道；
真正决定编码能否有效、系统能否可靠的，是信源和信道各自的统计特性。
\end{strand}

\block{通信技术的两个性能维度}
\trigger{当题目问“通信系统为什么要统计建模”时，先分清两个目标。}
\noindent\begin{tabularx}{\linewidth}{@{}l L@{}}
\toprule
{\color{indigo}目标} & \multicolumn{1}{l}{\color{indigo}依赖的统计特性}\\
\midrule
有效性 & 主要取决于信源统计特性。信源越有冗余，越有压缩空间。\\
可靠性 & 主要取决于信道统计特性。噪声、干扰和差错概率决定抗干扰设计。\\
\bottomrule
\end{tabularx}

\begin{definition}[信源]
信源是产生消息或信息的对象。它的输出通常不是确定值，而是具有
\answerin[3.4cm]{随机性} 的随机对象。
\end{definition}

\begin{definition}[信道]
信道是传递信号经过的通道。它会影响传输的可靠性，因此常用信道的
\answerin[3.4cm]{统计特性} 来描述噪声和差错。
\end{definition}

\begin{definition}[信源编码]
信源编码是在分析信源统计特性的基础上，用尽可能少的码率表示信源输出。
无失真编码要求接收端能够完全恢复原消息；限失真编码允许在给定失真约束内恢复
\answerin[3cm]{近似消息}。
\end{definition}

\begin{yourturn}
把“有效性”和“可靠性”各自对应的对象填出来：
\[
\text{有效性}\longleftrightarrow \answerin[3cm]{信源},
\qquad
\text{可靠性}\longleftrightarrow \answerin[3cm]{信道}.
\]
\workspace[2]
\end{yourturn}

\recall{信源编码为什么首先要研究信源统计特性？}

\section{信源分类与单符号信源}
\warmth{0}\quad\whisper{先看一次输出一个随机符号的信源，这是后面序列信源的基本积木。}

\begin{strand}
信源可以按输出的对象来分：如果每次输出一个符号，叫单符号信源；
如果连续输出一串符号，叫消息序列信源。也可以按取值集合分为离散信源和连续信源。
\end{strand}

\block{分类地图}
\noindent\begin{tabularx}{\linewidth}{@{}l L@{}}
\toprule
{\color{indigo}分类角度} & \multicolumn{1}{l}{\color{indigo}常见类型}\\
\midrule
按取值集合 & 离散信源，也可称数字信源；连续信源，也可称模拟信源。\\
按输出结构 & 单符号信源；消息序列信源，也就是随机过程形式的信源。\\
\bottomrule
\end{tabularx}

\begin{definition}[离散单符号信源]
设离散单符号信源为随机变量 \(X\)，其取值集合为
\[
\Xcal=\{x_1,x_2,\cdots,x_n\}.
\]
若每个符号 \(x_i\) 出现的概率为 \(\pmf(x_i)\)，则该信源的统计描述为
\[
\left(
\begin{array}{cccc}
x_1 & x_2 & \cdots & x_n\\
\pmf(x_1) & \pmf(x_2) & \cdots & \pmf(x_n)
\end{array}
\right),
\]
其中
\[
0\leq \pmf(x_i)\leq 1,\qquad
\sum_{i=1}^{n}\pmf(x_i)=\answerin[1.2cm]{1}.
\]
\end{definition}

\begin{definition}[连续单符号信源]
若信源输出 \(X\) 是连续随机变量，例如 \(X\in(a,b)\)，则通常用概率密度函数
\(\pmf(x)\) 描述。这里的 \(\pmf(x)\) 不是某一点的概率；当 \(dx\) 很小时，
\[
\pmf(x)\,dx \approx P\{x\leq X<x+dx\}.
\]
更一般地，
\[
P\{a\leq X\leq b\}=\int_a^b \pmf(x)\,dx.
\]
也就是说，\(\pmf(x)\) 描述的是单位长度小区间中的
\answerin[3.4cm]{概率密度}。
\end{definition}

\begin{remark}
离散情形用“概率质量”看每个符号，连续情形用“概率密度”看区间附近。
二者的共同点是：都在回答信源输出落在某些取值上的可能性。
本章聚焦离散信源；连续信源及其微分熵将在后续课程中展开。
\end{remark}

\begin{yourturn}
判断下面两个对象分别属于哪类信源：
\[
\text{二进制符号 } \{0,1\}:\ \answerin[3cm]{离散信源},
\qquad
\text{模拟电压 } X(t):\ \answerin[3cm]{连续信源}.
\]
\workspace[2]
\end{yourturn}

\recall{为什么连续信源中 \(\pmf(x)\) 通常叫概率密度而不是概率？}

\section{离散消息序列信源的统计描述}
\warmth{0}\quad\whisper{序列信源的关键变化：输出不再是一个符号，而是一整段符号向量。}

\begin{strand}
消息序列信源关心的不是单个输出符号，而是一段长度为 \(L\) 的符号序列。
因此统计描述从一维概率分布升级为 \(L\) 维联合概率分布。
\end{strand}

\block{从符号集合到序列集合}
设单个符号的取值集合为 \(\Xcal\)，则长度为 \(L\) 的序列取值集合可写成
\[
\Xcal^L=\Xcal\times \Xcal\times\cdots\times \Xcal,
\]
一次输出的消息序列为
\[
X^L=(X_1,X_2,\cdots,X_L),
\]
它是一个 \(L\) 维随机向量。某个具体序列可写为
\[
x^L=(x_1,x_2,\cdots,x_L).
\]

\begin{definition}[\(L\) 维联合分布]
离散消息序列信源的统计特性由联合概率
\[
\pmf(x^L)=\pmf(x_1,x_2,\cdots,x_L)
\]
描述。一般地，它可以按条件概率链式分解为
\[
\pmf(x_1,x_2,\cdots,x_L)
=\pmf(x_1)\pmf(x_2|x_1)\cdots
\pmf(x_L|x_1,x_2,\cdots,x_{L-1}).
\]
这个式子说明：序列中后一个符号的概率可能依赖于
\answerin[4cm]{前面已经出现的符号}。
\end{definition}

\begin{proposition}[序列取值个数]
若单符号信源有 \(n\) 个可能取值，则长度为 \(L\) 的离散消息序列共有
\[
n^L
\]
个可能序列。
\end{proposition}

\begin{example}
若单符号集合是 \(\{0,1\}\)，长度 \(L=3\)，则全部序列为
\[
000,\ 001,\ 010,\ 011,\ 100,\ 101,\ 110,\ 111.
\]
这里共有 \(2^3=8\) 个可能序列。
\end{example}

\begin{yourturn}
若四进制符号集合为 \(\{0,1,2,3\}\)，序列长度 \(L=2\)，可能序列数为
\[
\answerin[2cm]{16}.
\]
\workspace[2]
\end{yourturn}

\recall{为什么序列信源的概率模型一般要用联合分布，而不是只写每个符号自己的概率？}

\section{无记忆信源与等概二进制例子}
\warmth{0}\quad\whisper{无记忆是最常见的简化：过去的符号不影响现在的符号。}

\begin{strand}
真实序列信源可能有记忆，比如自然语言中前后字符相关。无记忆信源是一个重要特例：
序列中各符号统计独立，所以联合分布可以直接分解成乘积。
\end{strand}

\begin{definition}[无记忆离散信源]
若序列中的各个符号相互统计独立，则称该序列信源为无记忆信源。此时
\[
\pmf(x_1,x_2,\cdots,x_L)=
\prod_{\ell=1}^{L}\pmf(x_\ell).
\]
它的核心条件是序列中前后符号之间没有
\answerin[3.6cm]{统计依赖}。
\end{definition}

\begin{remark}
有记忆信源的相关性比较复杂，不能简单把联合概率写成各项乘积。
例如某些文字、语音或图像信源，当前符号往往受前面符号影响。
后面的有记忆信源编码会围绕一个目标展开：先解除或削弱相关性，再分别编码近似独立的分量。
\end{remark}

\begin{example}[三位等概二进制序列]
设单个符号 \(0\) 和 \(1\) 等概出现：
\[
\pmf(0)=\pmf(1)=\frac{1}{2}.
\]
当 \(L=3\) 时，可能序列共有 \(2^3=8\) 个。若信源无记忆，则每个具体序列的概率为
\[
\pmf(x_1,x_2,x_3)
=\pmf(x_1)\pmf(x_2)\pmf(x_3)
=\frac{1}{2}\cdot\frac{1}{2}\cdot\frac{1}{2}
=\frac{1}{8}.
\]
\end{example}

\begin{yourturn}
仍设 \(\pmf(0)=\pmf(1)=1/2\)。请写出序列 \(101\) 的概率计算：
\[
\pmf(101)=\pmf(1)\pmf(0)\pmf(1)
=\answerin[1.5cm]{\frac{1}{2}}\cdot\answerin[1.5cm]{\frac{1}{2}}\cdot\answerin[1.5cm]{\frac{1}{2}}
=\answerin[1.5cm]{\frac{1}{8}}.
\]
\workspace[2]
\end{yourturn}

\begin{yourturn}
如果 \(\pmf(0)=0.7,\ \pmf(1)=0.3\)，且仍为无记忆信源，请计算序列 \(001\) 的概率：
\[
\pmf(001)=\answerin[1.5cm]{0.7}\cdot\answerin[1.5cm]{0.7}\cdot\answerin[1.5cm]{0.3}
=\answerin[2cm]{0.147}.
\]
\workspace[2]
\end{yourturn}

\begin{strand}
复习时可以按这个顺序默写：信源产生消息，信源编码利用统计冗余；单符号信源写
\(X\) 和 \(\pmf(x_i)\)；序列信源写 \(X^L\) 和联合分布；无记忆时联合分布变成乘积。
接下来要问的是：这些概率分布对应的“不确定性”到底有多大，以及这种不确定性如何限制压缩码率？
\end{strand}

\section{信息熵：单符号离散信源的不确定性}
\warmth{0}\quad\whisper{熵不是某次收到的信息量，而是信源在统计意义上的平均不确定性。}

\begin{strand}
前面我们已经把信源写成概率分布。信息熵 \(H(X)\) 做的事，是把这个概率分布压缩成一个数：
信源平均有多不确定，平均需要多少信息才能消除这种不确定性。
\end{strand}

\block{定义}
\trigger{当题目给出离散信源 \(X\)、取值集合 \(\Xcal=\{x_1,\cdots,x_n\}\) 及概率 \(\pmf(x_i)\) 时，用这个公式。}

\begin{definition}[信息熵]
设离散单符号信源 \(X\) 的概率分布为 \(\pmf(x_i)\)。它的信息熵定义为
\[
H(X)=\mathbb{E}\big[-\log \pmf(X)\big]
=-\sum_{i=1}^{n}\pmf(x_i)\log \pmf(x_i).
\]
其中 \(\mathbb{E}\) 表示数学期望。信息熵度量的是信源的
\answerin[4cm]{平均不确定性}。
\end{definition}

\block{对数底与单位}
\noindent\begin{tabularx}{\linewidth}{@{}l l L@{}}
\toprule
{\color{indigo}对数底} & {\color{indigo}单位} & \multicolumn{1}{l}{\color{indigo}含义}\\
\midrule
\(2\) & bit & 最常用，表示二进制信息单位。\\
\(e\) & nat & 自然对数单位。\\
\(10\) & Hartley 或十进制单位 & 十进制对数单位。\\
\bottomrule
\end{tabularx}

\begin{remark}
若某个符号概率为 \(0\)，约定 \(0\log 0=0\)。直觉上，不可能发生的符号不会贡献平均不确定性。
\end{remark}

\block{二进制信源}
\begin{example}[二进制熵函数]
设二进制信源 \(X\in\{0,1\}\)，且
\[
\pmf(0)=p,\qquad \pmf(1)=1-p.
\]
若以 \(2\) 为对数底，则
\[
H(X)=-p\log_2 p-(1-p)\log_2(1-p)\triangleq h(p).
\]
当 \(p=\frac12\) 时，两个符号等概，二进制熵达到最大值 \(1\) bit。
\end{example}

\block{\(n\) 进制信源的最大熵}
\begin{proposition}[等概分布最大熵]
若 \(n\) 进制离散信源 \(X\) 的取值集合为 \(\Xcal=\{x_1,\cdots,x_n\}\)，则
\[
H(X)\leq \log n.
\]
等号成立当且仅当
\[
\pmf(x_1)=\pmf(x_2)=\cdots=\pmf(x_n)=\frac{1}{n}.
\]
\end{proposition}

\begin{proof}
记 \(p_i=\pmf(x_i)\)。若某些 \(p_i=0\)，这些项按约定
\(0\log 0=0\)，不会贡献熵。先考虑所有 \(p_i>0\) 的情形：
\[
H(X)=-\sum_{i=1}^{n}p_i\log p_i
=\sum_{i=1}^{n}p_i\log\frac{1}{p_i}.
\]
因为 \(\log t\) 是凹函数，由 Jensen 不等式，
\[
\sum_{i=1}^{n}p_i\log\frac{1}{p_i}
\leq
\log\left(\sum_{i=1}^{n}p_i\frac{1}{p_i}\right)
=\log n.
\]
因此 \(H(X)\leq \log n\)。

等号成立时，Jensen 不等式的等号条件要求
\[
\frac{1}{p_1}=\frac{1}{p_2}=\cdots=\frac{1}{p_n},
\]
也就是
\[
p_1=p_2=\cdots=p_n.
\]
再由 \(\sum_{i=1}^{n}p_i=1\)，得到
\[
p_1=p_2=\cdots=p_n=\frac{1}{n}.
\]
反过来，如果 \(p_i=1/n\)，则
\[
H(X)=-\sum_{i=1}^{n}\frac{1}{n}\log\frac{1}{n}
=\log n,
\]
所以等号确实成立。

若只有 \(k<n\) 个符号概率非零，则同样的证明给出
\[
H(X)\leq \log k<\log n,
\]
因此在 \(n\) 个可能符号中达到最大熵时，必须所有符号都以相同概率出现。
\end{proof}

\begin{yourturn}
复述最大熵证明的关键一步：把
\[
H(X)=\sum_{i=1}^{n}p_i\log\frac{1}{p_i}
\]
看成对凹函数 \(\log t\) 的加权平均，因此由 Jensen 不等式得到
\[
H(X)\leq
\log\left(\sum_{i=1}^{n}p_i\frac{1}{p_i}\right)
=\answerin[1.6cm]{\log n}.
\]
\workspace[2]
\end{yourturn}

\begin{strand}
\keyword{信息量}描述一次接收后不确定性减少了多少；\keyword{信息熵}描述接收前信源平均有多不确定。
前者偏向一次事件，后者偏向整个信源的统计平均。
对无失真离散信源编码，熵会成为平均码长的理论下界；对限失真编码，这个角色会由
\(R(D)\) 接过去。
\end{strand}

\recall{为什么等概信源的不确定性最大？}

\section{条件熵与联合熵}
\warmth{0}\quad\whisper{两个随机信源放在一起时，要区分“总不确定性”和“知道一个以后剩下多少”。}

\begin{strand}
现在有两个离散信源 \(X\) 和 \(Y\)。如果知道了 \(X\)，对 \(Y\) 的不确定性通常会下降。
这个剩余不确定性就是条件熵；把 \((X,Y)\) 当成一个整体看，就是联合熵。
\end{strand}

\block{条件熵}
\trigger{当题目出现“已知 \(X\) 后 \(Y\) 还剩多少不确定性”时，用 \(H(Y|X)\)。}

\begin{definition}[条件熵]
设 \(X\) 的取值为 \(x_1,\cdots,x_n\)，\(Y\) 的取值为 \(y_1,\cdots,y_m\)。条件熵定义为
\[
H(Y|X)
=-\sum_{i=1}^{n}\sum_{j=1}^{m}
\pmf(x_i,y_j)\log \pmf(y_j|x_i)
=\mathbb{E}\big[-\log \pmf(Y|X)\big].
\]
它度量的是已知 \(X\) 后，\(Y\) 的
\answerin[4cm]{剩余不确定性}。
\end{definition}

\begin{definition}[联合熵]
联合熵把二元随机向量 \((X,Y)\) 看成一个整体：
\[
H(X,Y)
=-\sum_{i=1}^{n}\sum_{j=1}^{m}
\pmf(x_i,y_j)\log \pmf(x_i,y_j).
\]
它度量两个信源共同输出时的总体不确定性。
\end{definition}

\begin{proposition}[熵的链式法则]
联合熵可以分解为
\[
H(X,Y)=H(X)+H(Y|X)
=H(Y)+H(X|Y).
\]
\end{proposition}

\begin{proof}
从联合概率的链式分解开始：
\[
\pmf(x_i,y_j)=\pmf(x_i)\pmf(y_j|x_i).
\]
代入联合熵定义，并把乘积的对数拆成和：
\[
-\sum_{i,j}\pmf(x_i,y_j)\log\big[\pmf(x_i)\pmf(y_j|x_i)\big]
\]
\[
=-\sum_{i,j}\pmf(x_i,y_j)\log \pmf(x_i)
-\sum_{i,j}\pmf(x_i,y_j)\log \pmf(y_j|x_i).
\]
第一项中 \(\sum_j\pmf(x_i,y_j)=\pmf(x_i)\)，所以它等于 \(H(X)\)；
第二项正是 \(H(Y|X)\)。因此
\[
H(X,Y)=H(X)+H(Y|X).
\]
同理也可以先写 \(\pmf(x_i,y_j)=\pmf(y_j)\pmf(x_i|y_j)\)，得到
\[
H(X,Y)=H(Y)+H(X|Y).
\]
\end{proof}

\begin{yourturn}
复述链式法则证明中最关键的一步：
\[
\log\big[\pmf(x_i)\pmf(y_j|x_i)\big]
=\answerin[4cm]{\log \pmf(x_i)+\log \pmf(y_j|x_i)}.
\]
\workspace[2]
\end{yourturn}

\block{Shannon 不等式}
\begin{proposition}[条件不会增加不确定性]
\[
H(Y|X)\leq H(Y),\qquad H(X|Y)\leq H(X).
\]
也就是说，知道另一个随机变量，只会减少或保持不确定性，不会增加平均不确定性。
\end{proposition}

\begin{yourturn}
若 \(X\) 和 \(Y\) 相互独立，则 \(Y\) 在已知 \(X\) 后没有变得更确定，因此
\[
H(Y|X)=\answerin[2cm]{H(Y)},
\qquad
H(X,Y)=\answerin[2cm]{H(X)}+\answerin[2cm]{H(Y)}.
\]
\workspace[2]
\end{yourturn}

\recall{为什么 \(H(X,Y)\) 既可以写成 \(H(X)+H(Y|X)\)，也可以写成 \(H(Y)+H(X|Y)\)？}

\section{互信息：不确定性的减少量}
\warmth{0}\quad\whisper{互信息把“收到一个信源后，对另一个信源多了解了多少”变成公式。}

\begin{strand}
如果 \(X\) 是信源，\(Y\) 是信宿或接收端观察到的量，那么互信息 \(\info(X;Y)\)
衡量的是：观察 \(Y\) 以后，关于 \(X\) 的不确定性平均减少了多少。
\end{strand}

\block{定义与等价写法}
\begin{definition}[互信息]
互信息定义为
\[
\info(X;Y)=H(X)-H(X|Y).
\]
等价地，也可以写成
\[
\info(X;Y)=H(Y)-H(Y|X)
=H(X)+H(Y)-H(X,Y).
\]
因此，互信息可以看成 \(X\) 与 \(Y\) 共享的那部分信息。
\end{definition}

\begin{proposition}[概率形式]
互信息也可以写成联合分布与边缘分布的比值：
\[
\info(X;Y)
=\sum_{i=1}^{n}\sum_{j=1}^{m}
\pmf(x_i,y_j)
\log\frac{\pmf(x_i,y_j)}{\pmf(x_i)\pmf(y_j)}.
\]
等价地，
\[
\info(X;Y)
=\mathbb{E}\left[\log\frac{\pmf(X|Y)}{\pmf(X)}\right]
=\mathbb{E}\left[\log\frac{\pmf(Y|X)}{\pmf(Y)}\right].
\]
\end{proposition}

\begin{proof}
从
\[
\info(X;Y)=H(X)-H(X|Y)
\]
出发。因为
\[
\pmf(x_i)=\sum_{j=1}^{m}\pmf(x_i,y_j),
\]
所以
\[
\begin{aligned}
H(X)
&=-\sum_{i=1}^{n}\pmf(x_i)\log \pmf(x_i)\\
&=-\sum_{i=1}^{n}\sum_{j=1}^{m}
\pmf(x_i,y_j)\log \pmf(x_i).
\end{aligned}
\]
另一方面，
\[
H(X|Y)
=-\sum_{i=1}^{n}\sum_{j=1}^{m}
\pmf(x_i,y_j)\log \pmf(x_i|y_j).
\]
两式相减，得到
\[
\begin{aligned}
\info(X;Y)
&=\sum_{i=1}^{n}\sum_{j=1}^{m}
\pmf(x_i,y_j)
\left[\log \pmf(x_i|y_j)-\log \pmf(x_i)\right]\\
&=\sum_{i=1}^{n}\sum_{j=1}^{m}
\pmf(x_i,y_j)
\log\frac{\pmf(x_i|y_j)}{\pmf(x_i)}.
\end{aligned}
\]
再由条件概率
\[
\pmf(x_i|y_j)=\frac{\pmf(x_i,y_j)}{\pmf(y_j)}
\]
可得
\[
\info(X;Y)
=\sum_{i=1}^{n}\sum_{j=1}^{m}
\pmf(x_i,y_j)
\log\frac{\pmf(x_i,y_j)}{\pmf(x_i)\pmf(y_j)}.
\]
这也说明
\[
\info(X;Y)
=\mathbb{E}\left[\log\frac{\pmf(X|Y)}{\pmf(X)}\right].
\]
同理从 \(\info(X;Y)=H(Y)-H(Y|X)\) 出发，可以得到
\[
\info(X;Y)
=\mathbb{E}\left[\log\frac{\pmf(Y|X)}{\pmf(Y)}\right].
\]
\end{proof}

\begin{yourturn}
复述互信息等价式证明的关键动作：
\[
H(X)=-\sum_i \pmf(x_i)\log\pmf(x_i)
=-\sum_{i,j}\answerin[2.8cm]{\pmf(x_i,y_j)}\log\pmf(x_i),
\]
所以 \(H(X)-H(X|Y)\) 可以合并成同一个 \((i,j)\) 双重求和。
\workspace[2]
\end{yourturn}

\block{性质}
\noindent\begin{tabularx}{\linewidth}{@{}l L@{}}
\toprule
{\color{indigo}性质} & \multicolumn{1}{l}{\color{indigo}含义}\\
\midrule
非负性 & \(\info(X;Y)\geq 0\)。平均来看，观察不会让不确定性“负减少”。\\
对称性 & \(\info(X;Y)=\info(Y;X)\)。\(Y\) 中关于 \(X\) 的信息量，等于 \(X\) 中关于 \(Y\) 的信息量。\\
独立时为零 & 若 \(X\) 与 \(Y\) 独立，则 \(\info(X;Y)=0\)。\\
完全确定时最大 & 若 \(Y\) 能完全确定 \(X\)，则 \(\info(X;Y)=H(X)\)。\\
\bottomrule
\end{tabularx}

\begin{remark}
在固定信道转移概率 \(\pmf(y|x)\) 时，\(\info(X;Y)\) 是输入分布 \(\pmf(x)\) 的凹函数；
在固定输入分布 \(\pmf(x)\) 时，\(\info(X;Y)\) 是信道转移概率 \(\pmf(y|x)\) 的凸函数。
这解释了两个后续优化问题：固定信道 \(\pmf(y|x)\) 时，最大化 \(\info(X;Y)\) 得信道容量；
固定失真约束时，最小化 \(\info(X;\hat X)\) 得率失真函数 \(R(D)\)。
\end{remark}

\begin{yourturn}
用熵的图像直觉填空：
\[
\info(X;Y)=H(X)-H(X|Y)
\]
表示观察 \(Y\) 以后，关于 \(X\) 的不确定性减少了
\answerin[3cm]{\info(X;Y)}。
\workspace[2]
\end{yourturn}

\begin{yourturn}
若 \(X\) 与 \(Y\) 独立，则 \(\pmf(x_i,y_j)=\pmf(x_i)\pmf(y_j)\)。代入概率形式可得
\[
\info(X;Y)=\sum_{i,j}\pmf(x_i,y_j)\log\answerin[1.5cm]{1}=\answerin[1.5cm]{0}.
\]
\workspace[2]
\end{yourturn}

\loose{下一步接信源编码：无失真看 \(H(X)\)，限失真看
\(R(D)=\min \info(X;\hat X)\)。}

\section{无失真离散信源编码}
\warmth{0}\quad\whisper{这一节把前面的熵变成编码结论：若允许很长的分组，平均码率或平均码长可以逼近 \(H(X)\)。}

\begin{strand}
无失真离散信源编码的核心问题是：把长度为 \(L\) 的信源序列 \(X^L\)
映射成码字，再由码字恢复出 \(\hat X^L\)。如果分组足够长，高概率序列集中在典型集里，
于是只需要给典型序列分配码字，所需码字数量大约是 \(2^{L H(X)}\)。
\end{strand}

\block{编码模型}
\trigger{当题目出现 \(X^L\)、码字集合 \(\mathcal{U}^k\)、编码映射和译码映射时，先画这条链。}

\begin{definition}[分组信源编码]
设离散无记忆信源的长度 \(L\) 输出序列取值于 \(\mathcal{X}^L\)，码符号集合为
\(\mathcal{U}\)，每个码字长度为 \(k\)，则码字集合为 \(\mathcal{U}^k\)。编码映射和译码映射分别为
\[
f:\mathcal{X}^L\to \mathcal{U}^k,
\qquad
g:\mathcal{U}^k\to \mathcal{X}^L.
\]
整个过程可写为
\[
X^L \xrightarrow{\ f\ } U^k \xrightarrow{\ g\ } \hat X^L,
\]
其中 \(\hat X^L\) 是译码恢复出来的
\answerin[3cm]{信源序列}。
\end{definition}

\begin{remark}
若 \(k\) 为固定值，就是等长编码；若不同信源序列允许使用不同长度的码字，就是变长编码。
本节先看分组等长编码的基本极限。
\end{remark}

\block{什么叫无失真}
\begin{definition}[严格无失真与渐近无失真]
严格无失真要求对每个输入序列都有
\[
\hat X^L=X^L.
\]
更常用的渐近无失真用错误概率描述：
\[
P_e^{(L)}=P\{\hat X^L\neq X^L\}.
\]
若当 \(L\to\infty\) 时存在一族编码方案使
\[
P_e^{(L)}\to 0,
\]
则称可以实现渐近无失真编码。
\end{definition}

\begin{yourturn}
把“严格无失真”和“渐近无失真”的差别填出来：
\[
\text{严格无失真：}\hat X^L=\answerin[2cm]{X^L},
\qquad
\text{渐近无失真：}P_e^{(L)}\to\answerin[1.4cm]{0}.
\]
\workspace[2]
\end{yourturn}

\block{码率}
\begin{definition}[编码速率]
若码符号集合 \(\mathcal{U}\) 的大小为 \(M\)，每个信源分组长度为 \(L\)，码字长度为 \(k\)，
则编码速率定义为
\[
R=\frac{k}{L}\log_2 M
\]
bit/source symbol。它表示平均每个信源符号需要多少二进制位来表示。
\end{definition}

\begin{example}
若用二进制码字，\(M=2\)，则
\[
R=\frac{k}{L}.
\]
若每个 \(L=4\) 的信源分组用 \(k=6\) 个二进制码元表示，则码率为
\[
R=\frac{6}{4}=1.5
\]
bit/source symbol。
\end{example}

\begin{yourturn}
若码符号集合大小为 \(M=4\)，长度 \(L=5\) 的信源序列被编码成长度 \(k=6\) 的码字，则
\[
R=\frac{k}{L}\log_2 M
=\frac{6}{5}\log_2 4
=\answerin[1.6cm]{2.4}.
\]
\workspace[2]
\end{yourturn}

\block{熵给出的可达极限}
\begin{theorem}[离散无记忆信源无失真编码定理]
设离散无记忆信源 \(X\) 的熵为 \(H(X)\)。对任意 \(\epsilon>0\)，若编码速率满足
\[
R=\frac{k}{L}\log_2 M>H(X)+\epsilon,
\]
则当 \(L\) 足够大时，存在一族分组编码使错误概率
\[
P_e^{(L)}<\delta
\]
其中 \(\delta>0\) 可以任意小。反过来，若
\[
R<H(X),
\]
则不可能实现渐近无失真编码。
\end{theorem}

\begin{remark}
这个定理的严格证明要用典型序列和大数定律，后面讲无失真信源编码定理时再完整证明。
这里先保留结论和直觉：高概率序列集中在典型集里，典型集大小约为
\(2^{L H(X)}\)，所以只要码字数能覆盖典型集，错误概率就可以随 \(L\) 增大而变小。
\end{remark}

\block{典型序列的角色}
\begin{definition}[典型集]
对长度为 \(L\) 的信源序列，把概率接近 \(2^{-L H(X)}\) 的高概率序列放入典型集，
记作 \(A_\epsilon^{(L)}\)。当 \(L\) 足够大时，
\[
P\{X^L\in A_\epsilon^{(L)}\}>1-\delta,
\]
并且典型集大小大约满足
\[
\left|A_\epsilon^{(L)}\right|\leq 2^{L(H(X)+\epsilon)}.
\]
\end{definition}

\begin{strand}
典型集解释了为什么“不是所有 \(n^L\) 个序列都同等重要”。虽然长度为 \(L\) 的序列总数是
\(n^L\)，但实际概率质量几乎都集中在典型集里。编码器只需要重点照顾典型序列，
非典型序列造成的错误概率可以随着 \(L\) 增大而变小。
\end{strand}

\begin{yourturn}
用典型集语言补完整个可达性思路：
\[
P\{X^L\in A_\epsilon^{(L)}\}>1-\delta,
\qquad
\left|A_\epsilon^{(L)}\right|\leq \answerin[3.4cm]{2^{L(H(X)+\epsilon)}}.
\]
因此只要
\[
M^k\geq \left|A_\epsilon^{(L)}\right|,
\]
就有足够多码字覆盖典型序列，也就是
\[
R=\frac{k}{L}\log_2 M>\answerin[2.6cm]{H(X)+\epsilon}.
\]
\workspace[3]
\end{yourturn}

\block{变长编码：平均码长逼近熵}
\begin{strand}
等长编码把每个长度为 \(L\) 的信源序列都塞进同样长的码字；变长编码则让高概率序列用短码字，
低概率序列用长码字。它关心的不是固定码字长度 \(k\)，而是平均码长 \(\bar K\)。
\end{strand}

\begin{definition}[平均码长]
设长度为 \(L\) 的信源序列 \(x^L\) 被编码成 \(m\) 元码字，码长为 \(K(x^L)\)。
平均码长定义为
\[
\bar K=\sum_{x^L\in\Xcal^L}\pmf(x^L)K(x^L).
\]
因此
\[
\frac{\bar K}{L}
\]
表示每个信源符号平均需要的 \(m\) 元码符号数。
\end{definition}

\begin{theorem}[变长编码定理]
设 \(X\) 为离散无记忆信源，熵为 \(H(X)\)，码符号集合大小为 \(m\)。
对长度为 \(L\) 的信源扩展进行唯一可译变长编码时，任意唯一可译码都满足
\[
\frac{\bar K}{L}\geq \frac{H(X)}{\log_2 m}.
\]
另一方面，存在一族 \(m\) 元即时码，使
\[
\frac{\bar K}{L}<\frac{H(X)}{\log_2 m}+\frac{1}{L}.
\]
也就是说，通过增大分组长度 \(L\)，变长编码的每符号平均码长可以任意逼近熵下界。
\end{theorem}

\begin{proof}
先看下界。任意唯一可译码的码长必须满足 Kraft 不等式，因此它诱导的平均码长不可能低于
信源扩展熵对应的 \(m\) 元码符号数：
\[
\bar K\geq \frac{H(X^L)}{\log_2 m}.
\]
由于信源无记忆，
\[
H(X^L)=L H(X),
\]
所以
\[
\frac{\bar K}{L}\geq \frac{H(X)}{\log_2 m}.
\]

再看可达性。对每个序列取 Shannon 码长
\[
K(x^L)=\left\lceil -\log_m \pmf(x^L)\right\rceil,
\]
则
\[
K(x^L)<-\log_m \pmf(x^L)+1.
\]
对 \(x^L\) 求平均，得到
\[
\bar K<\frac{H(X^L)}{\log_2 m}+1
=\frac{L H(X)}{\log_2 m}+1.
\]
两边除以 \(L\)，即得
\[
\frac{\bar K}{L}<\frac{H(X)}{\log_2 m}+\frac{1}{L}.
\]
\end{proof}

\begin{yourturn}
复述变长编码定理的关键夹逼：
\[
\frac{H(X)}{\log_2 m}
\leq
\frac{\bar K}{L}
<
\answerin[3.2cm]{\frac{H(X)}{\log_2 m}+\frac{1}{L}}.
\]
当 \(m=2\) 时，\(\log_2 m=1\)，所以
\[
H(X)\leq \frac{\bar K}{L}<\answerin[2cm]{H(X)+\frac{1}{L}}.
\]
\workspace[2]
\end{yourturn}

\begin{definition}[编码效率]
把平均码长换算成 bit/source symbol：
\[
R=\frac{\bar K}{L}\log_2 m.
\]
编码效率定义为
\[
\eta=\frac{H(X)}{R}.
\]
当 \(R\) 越接近 \(H(X)\)，编码效率越接近 \(1\)。
\end{definition}

\begin{yourturn}
若采用二进制变长编码，且分组长度 \(L\) 足够大，使 \(1/L<\epsilon\)，则定理给出
\[
H(X)\leq \frac{\bar K}{L}<H(X)+\answerin[1.4cm]{\epsilon}.
\]
这说明二进制变长编码的平均码长可以逼近
\answerin[2cm]{\(H(X)\)}。
\workspace[2]
\end{yourturn}

\recall{为什么无失真离散信源编码的极限码率是 \(H(X)\)，而不是 \(\log n\)？}

\section{信息率失真函数 \texorpdfstring{\(R(D)\)}{R(D)}}
\warmth{0}\quad\whisper{这一节把互信息变成压缩极限：如果允许一点失真，最少需要保留多少信息？}

\begin{strand}
有失真信源编码的核心问题是：在允许平均失真不超过 \(D\) 的条件下，寻找最优的
试验信道 \(\answerin[2.5cm]{p(\hat x|x)}\)，使得互信息
\(\info(X;\hat X)\) 最小。这个最小值就是 \(\answerin[1.5cm]{R(D)}\)。
\end{strand}

\block{从无失真到允许失真}
\trigger{当题目不再要求 \(\hat X=X\)，而是给出一个“平均误差不超过 \(D\)”时，就进入率失真框架。}

在数字通信系统中，信源编码器与译码器的整体作用可以抽象为
\[
X \xrightarrow{\quad\text{编码/译码}\quad} \hat X.
\]
无失真编码要求 \(\hat X=X\)；有失真编码只要求二者之间的
\(\answerin[2cm]{平均失真}\) 不超过某个允许值 \(D\)。

\begin{definition}[有失真编码的目标]
设信源 \(X\) 的字母表为 \(\Xcal=\{x_1,x_2,\cdots,x_n\}\)，重构符号字母表为
\(\hat{\Xcal}=\{\hat x_1,\hat x_2,\cdots,\hat x_m\}\)。
有失真信源编码的目标是：在
\[
\mathbb{E}\left[d(X,\hat X)\right] \leq \answerin[1.2cm]{D}
\]
的条件下，使编码所需的信息率尽可能 \(\answerin[1.5cm]{小}\)。
\end{definition}

\block{失真测度}
\begin{definition}[失真测度]
失真测度是一个非负函数
\[
d:\Xcal\times\hat{\Xcal}\to \mathbb{R}^{+},
\]
其中 \(d(x_i,\hat x_j)\) 表示把 \(x_i\) 重构为 \(\hat x_j\) 所产生的失真程度。
通常要求
\[
d(x_i,\hat x_j)\geq 0,\qquad d(x_i,x_i)=\answerin[1cm]{0}.
\]
\end{definition}

\begin{example}[汉明失真]
对于离散信源，汉明失真只关心“对”或“错”：
\[
d(x_i,\hat x_j)=
\begin{cases}
\answerin[1cm]{0}, & i=j,\\
\answerin[1cm]{1}, & i\neq j.
\end{cases}
\]
它的含义是：译码正确时失真为零，任何错误都计为同样大小的失真。
\end{example}

\begin{example}[平方误差失真]
对于连续信源，常用的失真测度是
\[
d(x,\hat x)=\answerin[2cm]{(x-\hat x)^2}.
\]
若 \(X,\hat X\) 为向量，则可推广为欧氏距离平方
\(\|\mathbf{x}-\hat{\mathbf{x}}\|^2\)。
\end{example}

\begin{yourturn}
汉明失真和平方误差失真的侧重点有何不同？
\[
\text{汉明失真：}\quad d\text{ 只与 }\answerin[2cm]{\text{是否相同}}\text{ 有关。}
\]
\[
\text{平方误差失真：}\quad d\text{ 还与 }\answerin[2cm]{\text{误差大小}}\text{ 有关。}
\]
\workspace[2]
\end{yourturn}


\block{平均失真与允许失真度}
设信源分布为 \(\pmf(x_i)\)，试验信道为 \(\pmf(\hat x_j|x_i)\)，则联合分布为
\[
\pmf(x_i,\hat x_j)=\pmf(x_i)\pmf(\hat x_j|x_i).
\]

\begin{definition}[平均失真]
平均失真定义为
\[
\bar d
=
\sum_{i=1}^{n}\sum_{j=1}^{m}
\pmf(x_i,\hat x_j)d(x_i,\hat x_j)
=
\mathbb{E}\left[d(X,\hat X)\right].
\]
\end{definition}

\begin{definition}[允许失真度]
允许失真度 \(D\) 是平均失真的上界。一个试验信道可行，当且仅当
\[
\bar d
=
\sum_{i=1}^{n}\sum_{j=1}^{m}
\pmf(x_i)\pmf(\hat x_j|x_i)d(x_i,\hat x_j)
\leq \answerin[1cm]{D}.
\]
\end{definition}

\begin{yourturn}
当 \(D\) 变小时，对重构质量的要求变 \(\answerin[1.5cm]{高}\)，
因此所需信息率通常会变 \(\answerin[1.5cm]{大}\)；
当 \(D\) 变大时，信息率可以变 \(\answerin[1.5cm]{小}\)。
\workspace[2]
\end{yourturn}

\block{试验信道}
\begin{definition}[试验信道]
在率失真理论中，信源编码与译码的整体过程抽象为条件概率分布
\(\pmf(\hat x_j|x_i)\)，称为 \textbf{试验信道}。
它不是一个物理信道，而是描述编码/译码造成的随机映射。
\end{definition}

所有满足平均失真约束的试验信道构成集合
\[
\mathcal{P}_D
=
\left\{
\pmf(\hat x|x):
\mathbb{E}\left[d(X,\hat X)\right]\leq D
\right\}.
\]
也就是说，\(\pmf(\hat x|x)\in \mathcal{P}_D\) 当且仅当
\[
\sum_{i=1}^{n}\sum_{j=1}^{m}
\pmf(x_i)\pmf(\hat x_j|x_i)d(x_i,\hat x_j)
\leq D.
\]

\block{互信息作为“必须保留的信息量”}
\begin{proposition}[互信息的概率形式]
原信源 \(X\) 与重构信源 \(\hat X\) 之间的互信息为
\[
\info(X;\hat X)
=
\sum_{i=1}^{n}\sum_{j=1}^{m}
\pmf(x_i)\pmf(\hat x_j|x_i)
\log
\frac{\pmf(\hat x_j|x_i)}{\pmf(\hat x_j)},
\]
其中
\[
\pmf(\hat x_j)=\sum_{i=1}^{n}\pmf(x_i)\pmf(\hat x_j|x_i).
\]
\end{proposition}

互信息衡量：重构信源 \(\hat X\) 中保留了多少关于原信源 \(X\) 的信息。
在有失真压缩中，\(\info(X;\hat X)\) 越小，意味着需要保留的信息越少，
相应编码所需的信息率越 \(\answerin[1cm]{低}\)。


\block{信息率失真函数的定义}
\begin{definition}[信息率失真函数]
信息率失真函数定义为
\[
R(D)
=
\min_{\pmf(\hat x|x)\in \mathcal{P}_D}
\info(X;\hat X).
\]
展开写即
\[
R(D)
=
\min_{\pmf(\hat x|x)}
\info(X;\hat X)
\]
满足约束
\[
\sum_{i=1}^{n}\sum_{j=1}^{m}
\pmf(x_i)\pmf(\hat x_j|x_i)d(x_i,\hat x_j)
\leq D,
\]
\[
\sum_{j=1}^{m}\pmf(\hat x_j|x_i)=1,
\quad i=1,2,\cdots,n,
\]
\[
\pmf(\hat x_j|x_i)\geq 0,
\quad i=1,2,\cdots,n,\ j=1,2,\cdots,m.
\]
\end{definition}

\begin{strand}
\(R(D)\) 的含义是：\textbf{在平均失真不超过 \(D\) 的条件下，编码所需的最小信息率}。
它是一个 \(\answerin[3cm]{\text{带约束的优化问题}}\)，优化变量是试验信道
\(\pmf(\hat x|x)\)，目标函数是互信息。
\end{strand}

\block{率失真函数的基本性质}

\begin{proposition}[非负性]
由于互信息非负，\(\info(X;\hat X)\geq 0\)，所以
\[
R(D)\geq \answerin[1cm]{0}.
\]
\end{proposition}

\begin{proposition}[单调非增性]
若 \(D_1<D_2\)，则
\[
R(D_1)\ \answerin[1.5cm]{\geq}\ R(D_2).
\]
直观含义：允许的失真越大，可选的试验信道越多，所需的最小信息率不会增加。
\end{proposition}

\begin{proposition}[凸性]
率失真函数 \(R(D)\) 是关于 \(D\) 的凸函数，即对 \(0\leq \lambda\leq 1\)，有
\[
R(\lambda D_1+(1-\lambda)D_2)
\leq
\lambda R(D_1)+(1-\lambda)R(D_2).
\]
\end{proposition}

\begin{proof}
设两个最优试验信道分别达到 \((D_1,R(D_1))\) 和 \((D_2,R(D_2))\)。
引入辅助变量 \(\theta\in\{1,2\}\)，以概率 \(\lambda\) 取 \(\theta=1\)、
以概率 \(1-\lambda\) 取 \(\theta=2\)。令 \(\theta\) 与 \(X\) 独立，
并在 \(\theta=k\) 时使用第 \(k\) 个试验信道 \(p_k(\hat x|x)\)。

这样得到的新条件分布为
\[
p(\hat x|x)=\lambda p_1(\hat x|x)+(1-\lambda)p_2(\hat x|x),
\]
其平均失真为
\[
\lambda D_1+(1-\lambda)D_2.
\]

下面估计其互信息。由数据处理不等式，
\[
\info(X;\hat X)\leq \info(X;\hat X,\theta).
\]
而
\[
\info(X;\hat X,\theta)=\info(X;\theta)+\info(X;\hat X|\theta).
\]
因为 \(\theta\) 与 \(X\) 独立，\(\info(X;\theta)=0\)，所以
\[
\info(X;\hat X|\theta)
=\lambda\info(X;\hat X_1)+(1-\lambda)\info(X;\hat X_2)
=\lambda R(D_1)+(1-\lambda)R(D_2).
\]
因此
\[
R(\lambda D_1+(1-\lambda)D_2)
\leq
\lambda R(D_1)+(1-\lambda)R(D_2).
\]
\end{proof}

\begin{yourturn}
凸性的关键动作是：用 \(\answerin[2cm]{\text{时间共享}}\) 把两个编码方案组合起来。
新方案的平均失真是两个方案对应量的 \(\answerin[2cm]{\text{加权平均}}\)，
而新方案的互信息 \(\info(X;\hat X)\) 是加权平均的 \(\answerin[2cm]{\text{上界}}\)。
\workspace[2]
\end{yourturn}

\begin{proposition}[连续性]
在通常条件下，\(R(D)\) 是关于 \(D\) 的连续函数。
\end{proposition}


\block{两个特殊点}

\begin{proposition}[零失真点]
当 \(D=0\) 时，有失真编码退化为无失真编码。

若信源为离散信源，且零失真等价于完全正确恢复，则
\[
R(0)=\info(X;X)=\answerin[2cm]{H(X)}.
\]

若信源为连续信源，在平方误差等常见失真度量下，随着 \(D\to 0\)，
通常有
\[
R(D)\to \answerin[2cm]{\infty}.
\]
\end{proposition}

\begin{proposition}[最大允许失真]
存在一个阈值 \(D_{\max}\)，使得当 \(D\geq D_{\max}\) 时，
可以不传输任何关于 \(X\) 的信息，而在接收端固定输出某个重构符号 \(\hat x\)。
此时 \(X\) 与 \(\hat X\) 独立，\(\info(X;\hat X)=\answerin[1cm]{0}\)，故
\[
R(D)=0,\qquad D\geq D_{\max}.
\]
其中
\[
D_{\max}
=
\min_{\hat x\in \hat{\Xcal}}
\sum_{i=1}^{n}\pmf(x_i)d(x_i,\hat x).
\]
\end{proposition}

\begin{yourturn}
补全 \(D_{\max}\) 的直观含义：在 \(\answerin[2cm]{\text{不发送信息}}\)
的情况下，接收端固定选择一个最优重构符号，
所能达到的最小平均失真为 \(\answerin[2.5cm]{D_{\max}}\)。
\workspace[2]
\end{yourturn}

\block{率失真曲线的整体图像}
综合以上性质，\(R(D)\) 的典型形状满足：
\[
R(0)=H(X),\qquad R(D)\text{ 单调非增且凸},\qquad R(D)=0\ (D\geq D_{\max}).
\]
它刻画了 \(\answerin[3cm]{\text{压缩率}}\) 与
\(\answerin[3cm]{\text{重构质量}}\) 之间的最优折中。

\begin{yourturn}
用一句话概括 \(R(D)\) 回答的问题：
\[
\text{若允许平均失真不超过 }D\text{，最少需要保留多少}
\answerin[2cm]{\text{比特}}\text{ 来描述信源？}
\]
\workspace[2]
\end{yourturn}

\block{本节小结}
\noindent\begin{tabularx}{\linewidth}{@{}l L@{}}
\toprule
{\color{indigo}概念} & \multicolumn{1}{l}{\color{indigo}作用}\\
\midrule
失真测度 \(d(x,\hat x)\) & 刻画原符号与重构符号之间的差异。\\
平均失真 \(\bar d=\mathbb{E}[d(X,\hat X)]\) & 描述整体重构误差。\\
允许失真度 \(D\) & 平均失真的上界。\\
试验信道 \(\pmf(\hat x|x)\) & 描述编码/译码整体造成的随机映射。\\
率失真函数 \(R(D)\) & 在失真约束下，编码所需的最小信息率。\\
\bottomrule
\end{tabularx}

\recall{为什么 \(R(D)\) 关于 \(D\) 是单调非增的？从可行集合
\(\mathcal{P}_D\) 的变化来解释。}

\loose{下一章可以接一个具体信源的 \(R(D)\) 计算，例如二元对称信源在汉明失真下的率失真函数。}

\section{连续信源的限失真编码}
\warmth{0}\quad\whisper{连续信源要先变成离散信源，才能进入前几节的编码框架。}

\begin{strand}
声音、图像等模拟信号在时间和幅度上都是连续的。要把它们纳入限失真编码，
必须经历\keyword{采样}、\keyword{量化}、\keyword{编码}三步：
采样把时间离散化，量化把幅度离散化，编码把符号比特化。
其中采样定理保证时间离散化可以不丢信息，而量化则不可避免地引入失真，
成为连续信源限失真编码的核心。
\end{strand}

\block{从模拟信号到比特流}
\trigger{当题目给出一个连续波形并要求压缩时，先画出这条链路。}

连续信源编码的标准流程可以写成
\[
x(t) \xrightarrow{\quad\text{采样}\quad} x[n]=x(nT_s)
\xrightarrow{\quad\text{量化}\quad} x_q[n]
\xrightarrow{\quad\text{编码}\quad} \text{比特流}.
\]

\begin{definition}[采样]
采样是以固定时间间隔 $T_s$ 从连续信号 $x(t)$ 中提取样本，得到离散时间序列
\[
x[n]=x(nT_s),\qquad n\in\mathbb{Z}.
\]
$T_s$ 称为 \answerin[2cm]{采样周期}，$f_s=1/T_s$ 称为 \answerin[2cm]{采样频率}。
这一步完成的是时间的 \answerin[2cm]{离散化}。
\end{definition}

\begin{definition}[量化]
量化是把每个采样值的连续幅度映射到有限个 \answerin[2cm]{量化电平} 之一：
\[
Q:\mathbb{R}\to\{y_1,y_2,\dots,y_M\}.
\]
这一步完成的是幅度的 \answerin[2cm]{离散化}，也是失真的主要来源。
\end{definition}

\begin{definition}[编码]
编码是把量化后的离散符号转换成 \answerin[2cm]{二进制码字}，以便存储或传输。
这一步完成的是 \answerin[2cm]{符号化}。
\end{definition}

\begin{yourturn}
把三步与它们的作用对应起来：
\[
\text{采样：}\answerin[2.5cm]{\text{时间离散化}},\quad
\text{量化：}\answerin[2.5cm]{\text{幅度离散化}},\quad
\text{编码：}\answerin[2.5cm]{\text{符号化}}.
\]
\workspace[2]
\end{yourturn}

\block{采样的频谱直觉}

采样在时域上相当于用冲激串乘以原信号，在频域上使原频谱以 $f_s$ 为周期重复。
如果重复后的频谱互不重叠，就可以用滤波器把原信号完整恢复出来；
如果重叠，就会产生 \keyword{混叠}，信息永久丢失。

\block{低通采样定理}
\trigger{当信号频谱集中在低频附近时，先用这个定理判断最小采样率。}

\begin{theorem}[低通采样定理，Shannon--Nyquist]
设实信号 $x(t)$ 是带限的，最高频率为 $W$（即频谱只在 $[-W,W]$ 内非零）。
则当采样频率满足
\[
f_s \geq \answerin[1.5cm]{2W}
\]
时，$x(t)$ 可由样本 $x(nT_s)$ 完全重建。
最小采样率 $2W$ 称为 \answerin[2.5cm]{Nyquist 速率}，对应采样周期 $1/(2W)$ 称为 \answerin[2.5cm]{Nyquist 间隔}。
\end{theorem}

\begin{example}
某低通信号最高频率 $W=4\,\text{kHz}$，则 Nyquist 速率为
\[
f_s \geq 2W = \answerin[2cm]{8\,\text{kHz}}.
\]
电话语音通常取 $8\,\text{kHz}$ 采样，正是基于这一结论。
\end{example}

\begin{yourturn}
若信号最高频率为 $W=10\,\text{kHz}$，则不产生混叠的最小采样率为
\[
f_{s,\min}=\answerin[2cm]{20\,\text{kHz}}.
\]
\workspace[2]
\end{yourturn}

\block{混叠失真}

\begin{definition}[混叠]
若采样频率不满足 $f_s \geq 2W$，频谱复制后会发生 \answerin[2cm]{重叠}。
此时不同频率成分在采样后变得不可区分，这种现象称为 \answerin[2cm]{混叠}（aliasing）。
混叠一旦发生，后续滤波无法恢复原始信号。
\end{definition}

\begin{yourturn}
判断下列情况是否会产生混叠：
\[
W=3\,\text{kHz},\ f_s=5\,\text{kHz}:\ \answerin[2cm]{\text{会}};\qquad
W=3\,\text{kHz},\ f_s=8\,\text{kHz}:\ \answerin[2cm]{\text{不会}}.
\]
\workspace[2]
\end{yourturn}

\block{带通采样定理}
\trigger{当信号是频带受限而不是基带受限时，可以用更低的采样率。}

\begin{proposition}[带通采样定理]
设实带通信号占据频段 $[f_L,f_H]$，带宽 $B=f_H-f_L$。令 $m=\lfloor f_H/B\rfloor$。
为保证采样后频谱不重叠，采样率应满足
\[
f_s \geq \frac{2f_H}{m}.
\]
特别地，若 $f_H$ 是 $B$ 的整数倍，则最低采样率可降到
\[
f_{s,\min}=\answerin[1.5cm]{2B},
\]
远低于低通采样要求的 $2f_H$。
\end{proposition}

\begin{example}
设 $f_L=100\,\text{kHz}$，$f_H=110\,\text{kHz}$，则 $B=10\,\text{kHz}$，$m=\lfloor 110/10\rfloor=11$。
带通最小采样率为
\[
f_{s,\min}=\frac{2\times 110\,\text{kHz}}{11}=\answerin[2cm]{20\,\text{kHz}},
\]
而低通采样要求 $2f_H=220\,\text{kHz}$。
\end{example}

\begin{yourturn}
设带通信号 $f_L=20\,\text{kHz}$，$f_H=30\,\text{kHz}$，则带宽 $B=\answerin[1.5cm]{10\,\text{kHz}}$，
$m=\answerin[1.5cm]{3}$，带通最小采样率为
\[
f_{s,\min}=\answerin[2cm]{20\,\text{kHz}}.
\]
\workspace[2]
\end{yourturn}

\block{重建}

\begin{proposition}[理想重建]
若采样满足相应的采样定理，则原信号可由样本通过理想滤波器恢复。

对低通信号：用增益为 $T_s$、截止频率为 $f_s/2$ 的
\answerin[2.8cm]{理想低通滤波器}，恢复公式为
\[
x(t)=\sum_{n=-\infty}^{\infty} x(nT_s)\,
\operatorname{sinc}\!\left(\frac{t-nT_s}{T_s}\right),
\]
其中 $\operatorname{sinc}(u)=\sin(\pi u)/(\pi u)$。

对带通信号：用中心频率位于信号频带中心的
\answerin[2.8cm]{理想带通滤波器} 恢复。
\end{proposition}

\begin{proof}[基带信号的理想重建（从卷积出发）]
只证低通情形；带通情形只需把低通滤波器换成覆盖信号频带的带通滤波器。

采样在时域上等价于用冲激串调制原信号，得到
\[
x_s(t)=x(t)\sum_{n=-\infty}^{\infty}\delta(t-nT_s)
      =\sum_{n=-\infty}^{\infty} x(nT_s)\,\delta(t-nT_s).
\]
它的频谱 $X_s(f)$ 是原频谱 $X(f)$ 以 $f_s$ 为周期的重复；低通采样定理保证这些副本互不重叠。

为恢复 $x(t)$，用一个理想低通滤波器保留基带并补偿采样引入的幅度缩放：
\[
H(f)=
\begin{cases}
T_s, & |f|\leq f_s/2,\\
0, & |f|>f_s/2.
\end{cases}
\]
其冲激响应为
\[
h(t)=\mathscr{F}^{-1}\{H(f)\}
    =T_s\int_{-f_s/2}^{f_s/2} e^{\mathrm{j}2\pi f t}\,\mathrm{d}f
    =\operatorname{sinc}(t/T_s).
\]

把采样序列 $x_s(t)$ 与 $h(t)$ 做卷积，即得重建信号
\[
\begin{aligned}
\hat{x}(t)
&=x_s(t)*h(t)\\
&=\int_{-\infty}^{\infty}
   \sum_{n=-\infty}^{\infty} x(nT_s)\,\delta(\tau-nT_s)\,
   h(t-\tau)\,\mathrm{d}\tau\\
&=\sum_{n=-\infty}^{\infty} x(nT_s)\,h(t-nT_s).
\end{aligned}
\]
代入 $h(t)=\operatorname{sinc}(t/T_s)$，就得到 sinc 插值公式
\[
x(t)=\sum_{n=-\infty}^{\infty} x(nT_s)\,
\operatorname{sinc}\!\left(\frac{t-nT_s}{T_s}\right).
\]
\end{proof}

\begin{yourturn}
复述理想重建证明的关键动作：
\begin{enumerate}[label=(\arabic*)]
\item 采样在时域写成冲激串调制：
      $x_s(t)=x(t)\sum_n\delta(t-nT_s)=\sum_n x(nT_s)\delta(t-nT_s)$；
\item 用增益为 \answerin[1cm]{$T_s$}、截止频率为 \answerin[1.5cm]{$f_s/2$} 的理想低通滤波器取出基带；
\item 求出滤波器冲激响应 $h(t)=\answerin[2.5cm]{\operatorname{sinc}(t/T_s)}$；
\item 将 $x_s(t)$ 与 $h(t)$ 做 \answerin[1.5cm]{卷积}，得到 sinc 插值公式。
\end{enumerate}
\workspace[2]
\end{yourturn}

\block{标量量化的基本模型}

\begin{definition}[标量量化]
标量量化是对采样序列样值进行量化。
它是一个从 $\mathbb{R}$ 到量化电平集合的映射：
\[
Q:\mathbb{R}\to\{y_1,y_2,\dots,y_M\}.
\]
具体地，每个实值样本 $x$ 被划分到某个区间，并用该区间对应的代表值表示：
\[
Q(x)=y_k,\quad \text{当 } x\in(t_{k-1},t_k],\ k=1,2,\dots,M.
\]
其中 $M$ 是 \answerin[1.5cm]{量化电平数}，$\{t_k\}$ 是 \answerin[2cm]{决策边界}，$\{y_k\}$ 是 \answerin[2cm]{量化电平}。
量化输出就是第 9 节的重构变量：
\[
\hat X=Q(X)\in\{y_1,\dots,y_M\}.
\]
\end{definition}

\block{量化噪声、功率与信噪比}

\begin{definition}[量化误差与量化噪声功率]
量化误差定义为
\[
e=x-Q(x).
\]
若把 $e$ 看成加性噪声，则量化噪声功率为
\[
N_q=\mathbb{E}\big[(X-Q(X))^2\big].
\]
\end{definition}

\begin{definition}[量化信噪比]
设输入信号功率为 $P_X=\mathbb{E}[X^2]$，则量化信噪比定义为
\[
\mathrm{SNR}=\dfrac{P_X}{N_q},
\]
常用分贝表示为
\[
\mathrm{SNR}_{\mathrm{dB}}=10\log_{10}\dfrac{P_X}{N_q}.
\]
\end{definition}

\begin{example}
若输入功率 $P_X=4$，量化噪声功率 $N_q=0.01$，则
\[
\mathrm{SNR}=400,\qquad
\mathrm{SNR}_{\mathrm{dB}}=26\,\text{dB}.
\]
\end{example}

\block{均匀量化器}
\trigger{最简单也最常用的量化器：所有区间一样大。}

\begin{definition}[均匀量化器]
若各量化区间长度相等，即
\[
\Delta=t_k-t_{k-1},\qquad k=1,2,\dots,M,
\]
则称为均匀量化器，$\Delta$ 称为 \answerin[1.5cm]{量化步长}。
\end{definition}

\begin{proposition}[均匀输入的量化噪声功率]
设输入 $X$ 在 $[-V,V]$ 上均匀分布，采用 $M$ 电平均匀量化器，则量化步长 $\Delta=2V/M$，量化噪声功率为
\[
N_q=\answerin[2.4cm]{\dfrac{\Delta^2}{12}}.
\]
\end{proposition}

\begin{example}
设 $X$ 在 $[-1,1]$ 均匀分布，用 $M=8$ 电平均匀量化，则 $\Delta=2/8=0.25$，
\[
N_q=\dfrac{(0.25)^2}{12}\approx\answerin[2.4cm]{5.21\times 10^{-3}}.
\]
信号功率 $P_X=1/3$，所以
\[
\mathrm{SNR}=\dfrac{1/3}{5.21\times 10^{-3}}\approx\answerin[2cm]{64}
\approx \answerin[2cm]{18.1\,\text{dB}}.
\]
\end{example}

\begin{yourturn}
把均匀量化 SNR 与量化比特数 $R=\log_2 M$ 的关系填出来：
\[
\mathrm{SNR}_{\mathrm{dB}}\approx 6.02R + C,
\]
其中 $C$ 是与 \answerin[3cm]{\text{输入信号统计特性}} 有关的常数。
对 $[-V,V]$ 均匀输入，$C=\answerin[2cm]{0\,\text{dB}}$。
\workspace[2]
\end{yourturn}

\block{最佳量化器}
\trigger{给定量化电平数，想让量化噪声最小，就要求解 Lloyd--Max 条件。}

\begin{proposition}[Lloyd--Max 最优条件]
给定量化电平数 $M$，使量化噪声功率 $N_q$ 最小的标量量化器满足：
\begin{enumerate}
\item 决策边界在两相邻量化电平中点：
\[
t_k=\frac{y_k+y_{k+1}}{2},\qquad k=1,2,\dots,M-1.
\]
\item 每个量化电平是对应区间上的条件期望：
\[
y_k=\mathbb{E}[X\mid t_{k-1}\leq X<t_k].
\]
\end{enumerate}
\end{proposition}

\begin{proof}
把 $N_q$ 对 $t_k$ 和 $y_k$ 分别求偏导并令其为零即可。
对边界 $t_k$ 求导，导数为零要求 $t_k$ 到 $y_k$ 与 $y_{k+1}$ 的距离相等，即中点条件。
对电平 $y_k$ 求导，导数为零要求 $y_k$ 是区间 $[t_{k-1},t_k]$ 上的条件均值，即质心条件。
这两个条件通常需要 \answertodo{用什么方法联合求解}{迭代}。
\end{proof}

\begin{yourturn}
对均匀输入，Lloyd--Max 条件退化为均匀量化器：因为条件期望正好落在区间
\answerin[2cm]{中点}，决策边界也正好在相邻电平 \answerin[2cm]{中点}。
\workspace[2]
\end{yourturn}

\begin{strand}
采样定理解决的是“时间离散化会不会丢信息”；
量化解决的是“幅度离散化能省多少比特、会引入多少失真”。
把二者合起来，就得到连续信源限失真编码的实际框架：先采样，再量化，最后对量化符号做无失真或限失真编码。
这正是第 9 节 $R(D)$ 的物理入口：量化器的设计直接决定了失真 $D$，而编码器在 $D$ 约束下逼近 $R(D)$。
\end{strand}

\loose{下一节将讨论矢量量化：把多个样本一起量化，利用样本间相关性进一步降低率失真。}

\section{矢量量化}
\warmth{0}\quad\whisper{标量量化一次只压一个样值；矢量量化把相邻样值打包，利用它们之间的相关性。}

\begin{strand}
矢量量化是对采样序列中每 $k$ 个样值的联合量化，属于多维量化。
它充分利用各样值间的 \answerin[2.5cm]{统计相关性}，因此在相同失真约束下，
编码效率 \answerin[2cm]{高于标量量化}。
\end{strand}

\block{基本原理}
\trigger{先把 $k$ 个样值看成向量，再把 $k$ 维空间切成若干子空间。}

\begin{definition}[矢量量化]
将模拟信号采样序列中每 $k$ 个样值分一组，得到 $k$ 维向量
\[
X=(x_1,x_2,\dots,x_k),
\]
其统计特性由联合概率密度 $\pmf(x_1,x_2,\dots,x_k)$ 描述。
将 $k$ 维空间分割成 $L$ 个子空间 $\{C_i\}_{i=1}^L$，满足
\[
\bigcup_{i=1}^L C_i=\mathbb{R}^k.
\]
每个子空间 $C_i$ 对应一个 \answerin[2cm]{重构矢量}（也称量化矢量或码字）
\[
y_i=(y_{i1},y_{i2},\dots,y_{ik})\in\mathbb{R}^k,\qquad i=1,2,\dots,L.
\]
若输入向量落在 $C_i$ 内，即 $X\in C_i$，则将其量化为码字 $y_i$（即重构矢量）：
\[
Q(X)=y_i.
\]
矢量量化可看作映射
\[
Q:\mathbb{R}^k\to\mathbb{R}^k.
\]
量化输出就是第 9 节的重构变量：
\[
\hat X=Q(X)\in\{y_1,\dots,y_L\}.
\]
\end{definition}

\noindent\begin{tabularx}{\linewidth}{@{}l l L@{}}
\toprule
{\color{indigo}量化方式} & {\color{indigo}输入/输出} & \multicolumn{1}{l}{\color{indigo}关键对象}\\
\midrule
标量量化 & $Q:\mathbb{R}\to\{y_k\}$ & 决策边界 $\{t_k\}$、量化电平 $\{y_k\}$\\
矢量量化 & $Q:\mathbb{R}^k\to\mathbb{R}^k$ & 子空间 $\{C_i\}$、码本 $\{y_i\}$\\
\bottomrule
\end{tabularx}

\begin{yourturn}
用一句话复述矢量量化的基本流程：
\begin{enumerate}[label=(\arabic*)]
\item 每 \answerin[1cm]{$k$} 个样值组成向量 $X$；
\item 把 $\mathbb{R}^k$ 分成 \answerin[1.5cm]{$L$} 个子空间 $\{C_i\}$；
\item 若 $X\in C_i$，则输出码字 \answerin[1.5cm]{$y_i$}。
\end{enumerate}
\workspace[2]
\end{yourturn}

\block{失真测度与平均失真}
\trigger{最常用的失真度量是均方误差。}

\begin{definition}[矢量量化的失真测度]
对输入向量 $X$ 与重构矢量 $y_i$，常用均方误差失真：
\[
d(X,y_i)=\sum_{j=1}^{k}(x_j-y_{ij})^2=\|X-y_i\|_2^2.
\]
\end{definition}

\begin{definition}[总平均失真]
矢量量化的总平均失真定义为
\[
D=\sum_{i=1}^{L}P(X\in C_i)\cdot\mathbb{E}\big[d(X,y_i)\mid X\in C_i\big],
\]
也可写成积分形式：
\[
D=\sum_{i=1}^{L}P(X\in C_i)\int_{X\in C_i}d(X,y_i)\,\pmf(x)\,dx,\qquad X\in\mathbb{R}^k.
\]
其中 $P(X\in C_i)=\answerin[2cm]{\displaystyle\int_{C_i}\pmf(x)\,dx}$ 是向量落入子空间 $C_i$ 的概率。
\end{definition}

\begin{example}
设 $k=2$，$L=2$，$P(X\in C_1)=0.6$，$P(X\in C_2)=0.4$。
若 $\mathbb{E}[d(X,y_1)\mid X\in C_1]=1.2$，
$\mathbb{E}[d(X,y_2)\mid X\in C_2]=2.5$，则
\[
D=0.6\times 1.2+0.4\times 2.5=\answerin[2cm]{1.72}.
\]
\end{example}

\block{最佳矢量量化的两个必要条件}
\trigger{给定码本大小 $L$，想让平均失真 $D$ 最小，就要同时满足最近邻与质心。}

\begin{proposition}[最佳矢量量化的两个必要条件]
给定码本大小 $L$，使总平均失真 $D$ 最小的矢量量化器满足：
\begin{enumerate}
\item \textbf{最近邻原则}（分割空间）：
\[
C_i=\{X:d(X,y_i)\le d(X,y_j),\ j\neq i\}
=\{X:\|X-y_i\|_2^2\le\|X-y_j\|_2^2,\ j\neq i\}.
\]
\item \textbf{质心原则}（码字设计）：
\[
y_i=\frac{\displaystyle\int_{X\in C_i}X\,\pmf(x)\,dx}
           {\displaystyle\int_{X\in C_i}\pmf(x)\,dx}.
\]
\end{enumerate}
\end{proposition}

\begin{remark}[LBG 算法]
上述两条件互相依赖：固定码本时最优划分是最近邻，固定划分时最优码字是质心。
Linde、Buzo 和 Gray（1980）据此给出交替迭代算法，常称为 \answerin[2cm]{LBG 算法}：
先给定初始码本，再反复用最近邻更新划分、用质心更新码字，直到平均失真收敛。
\end{remark}

\begin{yourturn}
对照第 10 节 Lloyd--Max：标量里决策边界在相邻电平的
\answerin[1.5cm]{中点}，矢量里子空间由
\answerin[2cm]{最近邻} 原则划分；
标量里量化电平是区间上的
\answerin[1.5cm]{条件期望}，矢量里码字是子空间上的
\answerin[2cm]{质心}。
\workspace[2]
\end{yourturn}

\recall{LBG 两条件：最近邻与质心。}

\begin{strand}
标量量化把每个样值独立离散化，忽略了相邻样值之间的相关性；
矢量量化把多个样值作为一个向量整体编码，在 $k$ 维空间里用 $L$ 个码字近似整个联合分布。
这正是第 9 节 $R(D)$ 的实际实现路径之一：码本大小 $L$ 决定每向量码率 $\log_2 L$，
折合到每样值约为
\[
R\approx\dfrac{1}{k}\log_2 L,
\]
而 LBG 条件则在给定 $L$ 时最小化平均失真 $D$。
\end{strand}

\loose{下一节接有记忆信源编码：先削弱相关性，再对近似独立分量编码（如 transform coding）。}


\end{document}
